\begin{exercice*}
    \begin{enumerate}
        \item Calcule, sur plusieurs exemples, la somme de quatre entiers consécutifs.

        \item Comment peut-on trouver le résultat juste en connaissant le premier entier ?

        \item Pour montrer que cette conjecture est toujours vraie, on désigne le premier des quatre entiers par la lettre $n$. Exprime alors les trois autres.

        \item Calcule alors la somme de ces quatre entiers et démontre ta conjecture.

        \item Que peux-tu dire de la somme de cinq entiers consécutifs ? Justifie.
    \end{enumerate}
\end{exercice*}
\begin{corrige}
    %\setcounter{partie}{0} % Pour s'assurer que le compteur de \partie est à zéro dans les corrigés
    %\phantom{rrr}    
    \dots
\end{corrige}

